\begin{resumo}

A preparação para concursos públicos e avaliações acadêmicas exige o domínio de estratégias de aprendizado eficazes e a manutenção do bem estar mental. Porém, os softwares disponíveis no mercado são fragmentados, impedindo o estudante de acessar um único aplicativo a fim de auxiliá-lo. Este trabalho traz a fundamentação da preparação para o desenvolvimento de um aplicativo mobile gratuito e \textit{open source} que unifica o Timer Pomodoro, o uso de repetição de flashcards, gamificação e check-ins de humor. A metodologia adotada combina principalmente o framework Lean Inception para o alinhamento da visão de produto e definição do MVP. Diante disso, foi feito o levantamento de requisitos a partir de personas, jornadas de usuário e brainstorming de funcionalidades além da definição da arquitetura e do banco de dados. O referencial teórico valida as funcionalidades propostas com base em evidências acadêmicas sobre gestão do tempo, repetição espaçada, gamificação e aprendizagem autorregulada. Como resultado deste trabalho, entrega-se uma proposta, teoricamente fundamentada e centrada no usuário, incluindo protótipo de alta fidelidade, modelagem de dados, diagramas de arquitetura e o planejamento para a fase de desenvolvimento.

 \vspace{\onelineskip}

 \noindent
 \textbf{Palavras-chave}: aplicativo mobile; gestão do tempo; repetição espaçada; gamificação; aprendizagem autorregulada.
\end{resumo}
